\section{Discussion and Future Work}

The framework proposed here is, at present, a design rather than a validated
system, and several questions remain open. The most pressing is empirical:
does evolutionary refinement of MicroSims actually produce measurable learning
gains over hand-authored content, and at what cost? The DGM and AlphaEvolve
results in their respective domains are encouraging, but education is a noisier
domain with smaller effect sizes and far greater measurement difficulty, so
direct transfer of their success rates should not be assumed.

Several avenues merit investigation. Transferable mutation operators---pedagogical
strategies such as prediction prompts that win in one concept and propagate to
others---could dramatically accelerate the search, echoing the general-purpose
improvements the DGM discovered. The graph-structural mutations of
Section~4.3, in which the system proposes new bridging concepts, represent a
more ambitious form of self-improvement that modifies the search space itself.
And the skill-level self-improvement of Section~4.4 raises the possibility of
compounding second-order gains, though it also compounds the safety concerns
and would require correspondingly stronger oversight.

The convergence of three trends makes this an opportune moment for such work:
foundation models capable of generating high-quality educational
content~\cite{mccreary2025claudeskills}, mature learning-analytics standards
in xAPI and cmi5 that provide the empirical grounding, and a body of
demonstrated self-improvement systems that provide a proven architectural
template. The central thesis of this paper is that intelligent textbooks
possess all the structural ingredients---modular content, instrumented
delivery, and a graph-organized search space---to become, in the DGM's phrase,
capable of gathering their own stepping stones along a path that unfolds toward
continuously improving instruction.
